% Book 4, Year 2, chapter 13 -- Processes and system calls.
% Provenance (never printed): sources/book4/BRIEFS.md, chapter 13.
\chapter{Processes and System Calls}\label{ch:b2:processes-and-system-calls}

At the end of this chapter, a shell of sixty lines of C runs this script:

\omoutput[1-10]{code/bachelor-2/13-processes-and-system-calls/outputs/minish.txt}{The chapter's shell,
\code{minish}, reading its first commands from a file.}\label{out:b2:processes-and-system-calls:hook}

Typing \code{echo hello from a child process} in a shell looks like one action. It is at least two programs: the
shell, which keeps running, and \code{echo}, which runs, prints and ends. Someone has to start \code{echo}, give it its
words, wait for it, and learn whether it succeeded. Neither program can do this alone: starting a program, giving it
memory, and writing to the screen are the work of the operating system's kernel. This chapter is about the narrow
door through which programs ask the kernel for things, and about the few requests that create and run processes.
Everything here is about Linux and other Unix-like systems; the tool \code{strace} used to watch the requests exists
on Linux only.

\begin{recall}
Book 2: a program is a file of instructions, a \omterm{def:b2:processes-and-system-calls:process}{process} is a program being run, with its own memory and state; the
operating system lets several processes share the processor. Book 3, ch.~8: a shell reads a command line, cuts it into
words, and runs the program the first word names with the others as \emph{command-line arguments}; a command ends
with an \emph{exit status}, 0 for success; \emph{environment variables} are name--value pairs passed to every
program. Book 3: the memory of a C \omterm{def:b2:processes-and-system-calls:process}{process} holds code, global data, a heap and a stack.
\cref{ch:b2:assembly-and-calling-convention}: \code{ecall} transfers control to the environment with a request
number in \code{a7}.
\end{recall}

\section{The system call}\label{sec:b2:processes-and-system-calls:syscall}

\begin{definition}[User mode, kernel mode]\label{def:b2:processes-and-system-calls:modes}
A processor runs in \emph{user mode}\index{user mode} or in \emph{kernel mode}\index{kernel mode}. In user mode, the
instructions that drive devices, change the memory map or stop the machine are refused, and only the memory the
\omterm{def:b2:processes-and-system-calls:process}{process} was given is reachable. In kernel mode everything is allowed. The kernel of the operating system runs in
kernel mode; every other program runs in user mode.
\end{definition}

\begin{definition}[System call]\label{def:b2:processes-and-system-calls:syscall}
A \emph{system call}\index{system call} is a request from a user-mode program to the kernel. The program puts a
request number and its arguments in registers and executes a special instruction (\code{ecall} on RISC-V,
\code{syscall} on x86-64), which switches the processor to \omterm{def:b2:processes-and-system-calls:modes}{kernel mode} and jumps to a fixed entry point in the kernel;
the kernel checks the request, does the work, puts the result in a register, and returns to \omterm{def:b2:processes-and-system-calls:modes}{user mode} just after the
instruction.
\end{definition}

The C library wraps each \omterm{def:b2:processes-and-system-calls:syscall}{system call} in a function of the same name: \code{write(fd, buffer, n)} asks the kernel to
write $n$ bytes to the open file numbered \code{fd}, and descriptor 1 is the standard output. \code{printf} is a
different kind of function, a \emph{library call}: it formats its text in the program's own memory, keeps it in a
buffer, and calls \code{write} later, when the buffer is full, at a newline on a terminal, or when the program ends.

\omcode{code/bachelor-2/13-processes-and-system-calls/listings/hello_write.c}{5}{10}{One line through the
library, one straight to the kernel}

\omoutput{code/bachelor-2/13-processes-and-system-calls/outputs/strace_hello.txt}{The program, then its
\code{write} and \code{exit_group} \omterm{def:b2:processes-and-system-calls:syscall}{system calls} as \code{strace} reports them.}\label{out:b2:processes-and-system-calls:strace}

The ``library'' line comes second. On a terminal \code{printf} would have flushed it at the newline; here the output
is not a terminal, so the line waits in the buffer until \code{main} returns, while the \code{write} went to the kernel at
once. \code{strace} runs a program and prints each of its \omterm{def:b2:processes-and-system-calls:syscall}{system calls} with its arguments and result: two
\code{write}s, of 14 and 15 bytes, and \code{exit_group}, the call that ends the \omterm{def:b2:processes-and-system-calls:process}{process}, which has no result since
it never returns.

\begin{omfigure}
\begin{tikzpicture}[font=\small, box/.style={draw, rounded corners=2pt, align=center, fill=omDef!10,
  minimum height=8mm, inner sep=3pt}, w/.style={-{Stealth[length=2mm]}}]
\fill[omThm!8] (-0.6,-0.8) rectangle (11.4,-3.0);
\node[anchor=west, font=\scriptsize\bfseries] at (-0.5,1.05) {user mode};
\node[anchor=west, font=\scriptsize\bfseries] at (-0.5,-2.75) {kernel mode};
\draw[dashed] (-0.6,-0.8) -- (11.4,-0.8);
\node[box] (p) at (1.2,0.3) {\code{main}};
\node[box] (lib) at (4.6,0.3) {C library\\ \code{write()}};
\node[box] (ent) at (4.6,-1.9) {entry point:\\ check the request};
\node[box] (fs) at (8.6,-1.9) {file layer:\\ descriptor 1};
\node[box] (dev) at (8.6,0.3) {terminal, file\\ or pipe};
\draw[w] (p) -- node[above, font=\scriptsize] {call} (lib);
\draw[w] ([xshift=4mm]lib.south) -- node[pos=0.25, right, font=\scriptsize, align=left] {number in a\\ register,
\code{syscall}} ([xshift=4mm]ent.north);
\draw[w] (ent) -- (fs);
\draw[w] (fs) -- node[pos=0.78, right, font=\scriptsize] {the bytes} (dev);
\draw[w, omThm] ([xshift=-4mm]ent.north) -- node[pos=0.75, left, font=\scriptsize, align=right] {result in a register,\\
back in user mode} ([xshift=-4mm]lib.south);
\end{tikzpicture}
\par\smallskip{\small The path of a \code{write}. The program calls an ordinary library function, which executes
the system-call instruction; the processor switches to \omterm{def:b2:processes-and-system-calls:modes}{kernel mode} at a fixed entry point, where the kernel checks
the descriptor and the buffer before doing the work. The result returns in a register, in \omterm{def:b2:processes-and-system-calls:modes}{user mode}.}
\end{omfigure}

The crossing is not free: the processor changes mode, saves and restores state, and the kernel checks every
argument. A measurement compares a call of an ordinary function with the cheapest \omterm{def:b2:processes-and-system-calls:syscall}{system call}, \code{getppid}, which
only reads a number the kernel keeps:

\omcode{code/bachelor-2/13-processes-and-system-calls/listings/syscall_cost.c}{25}{32}{Two million calls of each
kind}

\omoutput{code/bachelor-2/13-processes-and-system-calls/outputs/syscall_cost.txt}{One run on the author's laptop,
compiled with \code{-O2}.}\label{out:b2:processes-and-system-calls:cost}

A \omterm{def:b2:processes-and-system-calls:syscall}{system call} costs here about two hundred function calls. That is the reason for \code{printf}'s buffer: a program
that printed a thousand lines with one \code{write} each would make a thousand crossings instead of three
(\cref{exo:b2:processes-and-system-calls:11}).

\section{Creating a process: fork}\label{sec:b2:processes-and-system-calls:fork}

\begin{definition}[Process, process identifier]\label{def:b2:processes-and-system-calls:process}
A \emph{process}\index{process} is a running program together with the resources the kernel keeps for it: its memory,
its registers when it is not running, its open files, its working directory, its environment. The kernel names each
process by a \emph{process identifier}\index{process identifier} (PID), a positive integer unique among the
processes alive at the same time.
\end{definition}

\begin{definition}[Parent, child, process tree]\label{def:b2:processes-and-system-calls:tree}
A \omterm{def:b2:processes-and-system-calls:process}{process} that creates another is its \emph{parent}\index{parent process}, and the created one its
\emph{child}\index{child process}. Since every \omterm{def:b2:processes-and-system-calls:process}{process} but the first is created by another, the processes form a
\emph{process tree}\index{process tree}, rooted at the first \omterm{def:b2:processes-and-system-calls:process}{process} the kernel starts.
\end{definition}

% ledger: b2:fork-returns
\begin{definition}[Fork]\label{def:b2:processes-and-system-calls:fork}
A \emph{fork}\index{fork} creates a new \omterm{def:b2:processes-and-system-calls:process}{process} by duplicating the caller. The \omterm{def:b2:processes-and-system-calls:syscall}{system call} \code{fork()} makes a
\omterm{def:b2:processes-and-system-calls:tree}{child} whose memory is a copy of the \omterm{def:b2:processes-and-system-calls:tree}{parent}'s, with the same open files and working directory, and both processes
continue from the return of \code{fork}. It returns 0 in the \omterm{def:b2:processes-and-system-calls:tree}{child} and the \omterm{def:b2:processes-and-system-calls:tree}{child}'s PID in the \omterm{def:b2:processes-and-system-calls:tree}{parent}, or $-1$ in the
\omterm{def:b2:processes-and-system-calls:tree}{parent}, with no \omterm{def:b2:processes-and-system-calls:tree}{child}, if it fails.
\end{definition}

One call, two returns: the only way for the code to know which \omterm{def:b2:processes-and-system-calls:process}{process} it is running in is the value \code{fork}
returned. The copy is real: a variable changed in the \omterm{def:b2:processes-and-system-calls:tree}{child} keeps its old value in the \omterm{def:b2:processes-and-system-calls:tree}{parent}.

\omcode{code/bachelor-2/13-processes-and-system-calls/listings/fork_basic.c}{8}{24}{Two processes, one program}

\omoutput{code/bachelor-2/13-processes-and-system-calls/outputs/fork_basic.txt}{Each \omterm{def:b2:processes-and-system-calls:process}{process} learns its role from
the value of \code{fork}.}

The program prints relations, never the numbers themselves, because PIDs change from one run to the next.
\code{waitpid}, the subject of \cref{sec:b2:processes-and-system-calls:wait}, makes the \omterm{def:b2:processes-and-system-calls:tree}{parent} wait until the
\omterm{def:b2:processes-and-system-calls:tree}{child} has ended, which fixes the order of the two lines. Linux does not copy the memory at the moment of the \omterm{def:b2:processes-and-system-calls:fork}{fork}:
the two processes share the pages, marked read-only, and a page is copied only when one of them writes to it. The
program cannot tell the difference, and a \omterm{def:b2:processes-and-system-calls:fork}{fork} followed at once by an \code{exec} copies almost nothing.

\begin{omfigure}
\begin{tikzpicture}[font=\small, mem/.style={draw, minimum width=3.6cm, minimum height=6mm, fill=omDef!8},
  w/.style={-{Stealth[length=2mm]}}]
\node[mem] (pa) at (0,1.3) {code};
\node[mem] (pb) at (0,0.7) {\code{books = 10}};
\node[mem] (pc) at (0,0.1) {stack: \code{r} $=$ child's PID};
\node[mem] (ca) at (6.5,1.3) {code};
\node[mem] (cb) at (6.5,0.7) {\code{books = 20}};
\node[mem] (cc) at (6.5,0.1) {stack: \code{r} $= 0$};
\node[above=2pt of pa, font=\small\bfseries] {parent};
\node[above=2pt of ca, font=\small\bfseries] {child};
\draw[w, dashed] (pa.east) -- node[above, font=\scriptsize] {copied by \code{fork}} (ca.west);
\draw[w, dashed] (pb.east) -- (cb.west);
\draw[w, dashed] (pc.east) -- (cc.west);
\node[font=\scriptsize, align=center] at (3.25,-0.75) {after the copy, each process changes only its own memory};
\end{tikzpicture}
\par\smallskip{\small A \omterm{def:b2:processes-and-system-calls:fork}{fork} of \code{fork_basic.c}. The two address spaces start equal, except for the value
\code{fork} returned; the \omterm{def:b2:processes-and-system-calls:tree}{child}'s assignment \code{books = 20} does not reach the \omterm{def:b2:processes-and-system-calls:tree}{parent}.}
\end{omfigure}

A \omterm{def:b2:processes-and-system-calls:process}{process} that calls \code{fork} several times, with its children doing the same, grows a tree. In
\code{fork_tree.c} each \omterm{def:b2:processes-and-system-calls:process}{process} names itself \code{p} followed by the indices of the \omterm{def:b2:processes-and-system-calls:fork}{forks} it came out of as the
\omterm{def:b2:processes-and-system-calls:tree}{child}, and the \omterm{def:b2:processes-and-system-calls:tree}{parent} waits for each \omterm{def:b2:processes-and-system-calls:tree}{child}, so that each subtree finishes before the next \omterm{def:b2:processes-and-system-calls:fork}{fork}:

\omcode{code/bachelor-2/13-processes-and-system-calls/listings/fork_tree.c}{7}{23}{Three \omterm{def:b2:processes-and-system-calls:fork}{forks} in a
row}

\omoutput{code/bachelor-2/13-processes-and-system-calls/outputs/fork_tree.txt}{Eight processes; each prints after
its children.}\label{out:b2:processes-and-system-calls:tree}

\begin{omfigure}
\begin{tikzpicture}[font=\small, level distance=11mm, sibling distance=24mm,
  every node/.style={draw, rounded corners=2pt, fill=omDef!10, inner sep=2pt, minimum width=9mm},
  level 2/.style={sibling distance=13mm}, level 3/.style={sibling distance=10mm}]
\node {\code{p}}
  child {node {\code{p0}}
    child {node {\code{p01}} child {node {\code{p012}}}}
    child {node {\code{p02}}}}
  child {node {\code{p1}} child {node {\code{p12}}}}
  child {node {\code{p2}}};
\end{tikzpicture}
\par\smallskip{\small The \omterm{def:b2:processes-and-system-calls:tree}{process tree} of \code{fork_tree.c}: an edge goes from a \omterm{def:b2:processes-and-system-calls:tree}{parent} to a \omterm{def:b2:processes-and-system-calls:tree}{child}. The \omterm{def:b2:processes-and-system-calls:fork}{fork} of
index $i$ is made by every \omterm{def:b2:processes-and-system-calls:process}{process} alive at that moment, and doubles their number.}
\end{omfigure}

\begin{proposition}[Counting forks]\label{prop:b2:processes-and-system-calls:fork-count}
If a \omterm{def:b2:processes-and-system-calls:process}{process} executes $n$ successive \code{fork} calls, and every \omterm{def:b2:processes-and-system-calls:process}{process} created continues from the same point and
executes the remaining calls, then $2^n$ processes exist at the end, counting the first.
\end{proposition}

The proof is \cref{exo:b2:processes-and-system-calls:8}. With $n = 3$ the output has eight lines, and the depth of
each \omterm{def:b2:processes-and-system-calls:process}{process}, the number of \omterm{def:b2:processes-and-system-calls:fork}{forks} it came out of as the \omterm{def:b2:processes-and-system-calls:tree}{child}, takes the value $d$ for $\binom{3}{d}$ of them.

\section{Running a program: exec}\label{sec:b2:processes-and-system-calls:exec}

A \omterm{def:b2:processes-and-system-calls:fork}{fork} makes a copy of the same program. To run another one, a \omterm{def:b2:processes-and-system-calls:process}{process} replaces its own program.

% ledger: b2:exec-no-return
\begin{definition}[Exec, argument vector]\label{def:b2:processes-and-system-calls:exec}
An \emph{exec}\index{exec} replaces the program of the calling \omterm{def:b2:processes-and-system-calls:process}{process} by a new one, read from a file: the code,
data, heap and stack are discarded and rebuilt, and the new program starts at its beginning, in the same \omterm{def:b2:processes-and-system-calls:process}{process},
with the same PID, open files and working directory. Its arguments are given as an \emph{argument vector}\index{argument
vector}, an array of strings ending with a null pointer, whose element 0 is by convention the program's name; it
becomes \code{argv} in the new program's \code{main}. A successful exec does not return; if it returns, it failed.
\end{definition}

The C library offers several front ends to the one \omterm{def:b2:processes-and-system-calls:syscall}{system call} \code{execve(path, argv, envp)}: \code{execvp} searches
the directories of the \code{PATH} variable for the program and keeps the environment; \code{execlp} takes the
arguments one by one; \code{execve} itself takes a full path and a new environment.

\omcode{code/bachelor-2/13-processes-and-system-calls/listings/exec_demo.c}{8}{28}{Three \omterm{def:b2:processes-and-system-calls:exec}{execs}: in two children,
then in the \omterm{def:b2:processes-and-system-calls:process}{process} itself}

\omoutput{code/bachelor-2/13-processes-and-system-calls/outputs/exec_demo.txt}{Each line after the first comes from
another program.}

\code{echo} receives \code{"two words"} as one argument: no shell cut it, the vector is passed as it is. \code{env},
started with an environment of two variables, prints exactly those. The last \omterm{def:b2:processes-and-system-calls:exec}{exec} replaces \code{exec_demo} itself
by \code{printf}, so the line after it is never printed, and the exit status of the \omterm{def:b2:processes-and-system-calls:process}{process} is \code{printf}'s. The
\code{fflush} at the top matters for the same reason: an \omterm{def:b2:processes-and-system-calls:exec}{exec} discards the memory of the old program, buffer
included, and a line still waiting in \code{printf}'s buffer would be lost.

\section{Waiting for a child}\label{sec:b2:processes-and-system-calls:wait}

% ledger: b2:waitpid-macros
\begin{definition}[Waiting for a child]\label{def:b2:processes-and-system-calls:wait}
A \omterm{def:b2:processes-and-system-calls:process}{process} \emph{waits for a child}\index{waiting for a child} with \code{waitpid(pid, &status, 0)}, which blocks
until that \omterm{def:b2:processes-and-system-calls:tree}{child} has ended and returns its PID; \code{wait(&status)} does the same for any \omterm{def:b2:processes-and-system-calls:tree}{child}. The integer
\code{status} encodes how the \omterm{def:b2:processes-and-system-calls:tree}{child} ended: \code{WIFEXITED(status)} is true when it ended by returning from
\code{main} or calling \code{exit}, and \code{WEXITSTATUS(status)} is then the low 8 bits of its exit status;
\code{WIFSIGNALED(status)} is true when a signal killed it, and \code{WTERMSIG(status)} is the signal's number.
\end{definition}

\omcode{code/bachelor-2/13-processes-and-system-calls/listings/wait_status.c}{10}{17}{Reading a status}

\omoutput{code/bachelor-2/13-processes-and-system-calls/outputs/wait_status.txt}{Four children, four
endings.}\label{out:b2:processes-and-system-calls:status}

Signals are the subject of \cref{ch:b2:files-pipes-signals}; here it is enough that \code{raise(SIGTERM)} sends the
\omterm{def:b2:processes-and-system-calls:process}{process} the ``terminate'' signal, number 15, whose default effect is to end it. \code{exit(300)} gives 44, the low
8 bits of 300: an exit status is a number from 0 to 255. Shells report a death by signal $n$ as the status
$128 + n$, so that one number covers both cases.

Between its end and its \omterm{def:b2:processes-and-system-calls:tree}{parent}'s wait, a \omterm{def:b2:processes-and-system-calls:tree}{child} is in a peculiar state: its program, memory and files are gone, but
the kernel must keep its status until the \omterm{def:b2:processes-and-system-calls:tree}{parent} asks for it.

% ledger: b2:proc-state-z
\begin{definition}[Zombie, orphan]\label{def:b2:processes-and-system-calls:zombie}
A \emph{zombie}\index{zombie process} is a \omterm{def:b2:processes-and-system-calls:process}{process} that has ended but whose \omterm{def:b2:processes-and-system-calls:tree}{parent} has not yet waited for it: only
its entry in the kernel's \omterm{def:b2:processes-and-system-calls:process}{process} table remains, with its PID and its status. An \emph{orphan}\index{orphan
process} is a \omterm{def:b2:processes-and-system-calls:process}{process} whose \omterm{def:b2:processes-and-system-calls:tree}{parent} has ended before it; the kernel gives it a new \omterm{def:b2:processes-and-system-calls:tree}{parent}, the first \omterm{def:b2:processes-and-system-calls:process}{process} or a
\omterm{def:b2:processes-and-system-calls:process}{process} that asked to adopt orphans, which waits for it when it ends.
\end{definition}

On Linux the kernel shows each \omterm{def:b2:processes-and-system-calls:process}{process} as a directory \code{/proc/<pid>}, whose file \code{status} has a line
\code{State:}, with \code{Z} for a \omterm{def:b2:processes-and-system-calls:zombie}{zombie}. \code{zombie.c} reads that letter for a \omterm{def:b2:processes-and-system-calls:tree}{child} that has ended, before and
after waiting for it, then makes an \omterm{def:b2:processes-and-system-calls:zombie}{orphan}: a middle \omterm{def:b2:processes-and-system-calls:process}{process} starts a grandchild and ends at once, and the grandchild
watches its own \omterm{def:b2:processes-and-system-calls:tree}{parent} change.

\omcode{code/bachelor-2/13-processes-and-system-calls/listings/zombie.c}{30}{36}{A \omterm{def:b2:processes-and-system-calls:tree}{child} ended, not yet waited
for}

\omoutput{code/bachelor-2/13-processes-and-system-calls/outputs/zombie.txt}{A \omterm{def:b2:processes-and-system-calls:zombie}{zombie}, then nothing; an \omterm{def:b2:processes-and-system-calls:zombie}{orphan},
adopted.}

After the wait, \code{/proc} has no entry any more: the letter \code{-} is the program's way of saying so. A
long-running \omterm{def:b2:processes-and-system-calls:tree}{parent} that never waits accumulates \omterm{def:b2:processes-and-system-calls:zombie}{zombies}, each holding a PID, and PIDs are a finite resource
(\cref{exo:b2:processes-and-system-calls:2}).

\begin{omfigure}
\begin{tikzpicture}[font=\small, node distance=2.6cm, every state/.style={minimum size=13mm, fill=omDef!10,
  inner sep=1pt}, >={Stealth[length=2mm]}]
\node[state, initial, initial text={\code{fork}}] (ready) {ready};
\node[state, right=of ready] (run) {running};
\node[state, below=1.6cm of run] (blk) {blocked};
\node[state, right=of run, fill=omThm!15] (z) {zombie};
\node[right=1.6cm of z, font=\small] (gone) {gone};
\path[->] (ready) edge[bend left] node[above, font=\scriptsize] {scheduled} (run)
  (run) edge[bend left] node[below, font=\scriptsize] {preempted} (ready)
  (run) edge node[right, font=\scriptsize] {waits (\code{read}, \code{waitpid})} (blk)
  (blk) edge[bend left=20] node[below left, font=\scriptsize] {event} (ready)
  (run) edge node[above, font=\scriptsize] {\code{exit}, signal} (z)
  (z) edge node[above, font=\scriptsize, align=center] {parent\\ waits} (gone);
\end{tikzpicture}
\par\smallskip{\small The states of a \omterm{def:b2:processes-and-system-calls:process}{process}, with the \omterm{def:b2:processes-and-system-calls:zombie}{zombie} state. Book 2's ready, running and blocked states
remain; ending leads to the \omterm{def:b2:processes-and-system-calls:zombie}{zombie} state, and only the \omterm{def:b2:processes-and-system-calls:tree}{parent}'s wait removes the \omterm{def:b2:processes-and-system-calls:process}{process}.}
\end{omfigure}

\section{A minimal shell}\label{sec:b2:processes-and-system-calls:shell}

The four calls are all a shell needs. Its main loop:

\begin{method}[The shell's loop]\label{met:b2:processes-and-system-calls:shell-loop}
Repeat until the input ends or the command is \code{exit}:
\begin{enumerate}
\item print a prompt and read one line;
\item cut it into words; an empty line is skipped;
\item if the first word names a \omterm{def:b2:processes-and-system-calls:builtin}{built-in command}, run it in the shell itself;
\item otherwise \omterm{def:b2:processes-and-system-calls:fork}{fork}; in the \omterm{def:b2:processes-and-system-calls:tree}{child}, \omterm{def:b2:processes-and-system-calls:exec}{exec} the program the first word names with the words as its \omterm{def:b2:processes-and-system-calls:exec}{argument vector}, and
if the \omterm{def:b2:processes-and-system-calls:exec}{exec} fails, report it and end the \omterm{def:b2:processes-and-system-calls:tree}{child} with status 127;
\item in the \omterm{def:b2:processes-and-system-calls:tree}{parent}, wait for the \omterm{def:b2:processes-and-system-calls:tree}{child}, and keep its status for the next \code{$?}.
\end{enumerate}
\end{method}

\begin{omfigure}
\begin{tikzpicture}[font=\small, node distance=7mm and 9mm, b/.style={draw, rounded corners=2pt, fill=omDef!10,
  align=center, minimum height=7mm, inner sep=3pt}, d/.style={draw, diamond, aspect=2.2, fill=omProp!12, inner
  sep=1pt, align=center}, w/.style={-{Stealth[length=2mm]}}]
\node[b] (read) {prompt, read a line};
\node[b, right=of read] (split) {split into words};
\node[d, right=of split] (bi) {built-in?};
\node[b, below=of bi] (fork) {\code{fork}};
\node[b, below left=7mm and 4mm of fork] (child) {child: \code{execvp}\\ or \code{_exit(127)}};
\node[b, below right=7mm and 4mm of fork] (parent) {parent: \code{waitpid},\\ keep the status};
\node[b, above=of bi] (run) {run it in the shell\\ (\code{cd}, \code{exit})};
\draw[w] (read) -- (split); \draw[w] (split) -- (bi);
\draw[w] (bi) -- node[right, font=\scriptsize] {no} (fork);
\draw[w] (bi) -- node[right, font=\scriptsize] {yes} (run);
\draw[w] (fork) -- (child); \draw[w] (fork) -- (parent);
\draw[w] (parent.south) -- ++(0,-0.4) -| (read.south);
\draw[w] (run.west) -| (read.north);
\end{tikzpicture}
\par\smallskip{\small The loop of \cref{met:b2:processes-and-system-calls:shell-loop}. Only the \omterm{def:b2:processes-and-system-calls:tree}{parent} goes
round again; the \omterm{def:b2:processes-and-system-calls:tree}{child} either becomes the command or ends.}
\end{omfigure}

The splitting is a separate file, so that it can be tested alone. It cuts the line in place, at blanks, and points
each word of the vector into the line; the word \code{$?} becomes the last status.

\omcode{code/bachelor-2/13-processes-and-system-calls/listings/minish_parse.c}{7}{20}{Cutting a line into
words}\label{lst:b2:processes-and-system-calls:split}

\omcode{code/bachelor-2/13-processes-and-system-calls/listings/minish.c}{11}{27}{An external command: \omterm{def:b2:processes-and-system-calls:fork}{fork}, \omterm{def:b2:processes-and-system-calls:exec}{exec},
wait}\label{lst:b2:processes-and-system-calls:run}

\omcode{code/bachelor-2/13-processes-and-system-calls/listings/minish.c}{38}{58}{The loop}

The \code{_exit(127)} after a failed \omterm{def:b2:processes-and-system-calls:exec}{exec} is essential. Without it the \omterm{def:b2:processes-and-system-calls:tree}{child} returns from \code{run} into the shell's
loop, and two shells read the same input from then on (\cref{exo:b2:processes-and-system-calls:6}). It is
\code{_exit}, not \code{exit}, so that the \omterm{def:b2:processes-and-system-calls:tree}{child} does not flush a copy of the \omterm{def:b2:processes-and-system-calls:tree}{parent}'s buffers. Here is the end of
the hook's script:

\omoutput[11-24]{code/bachelor-2/13-processes-and-system-calls/outputs/minish.txt}{The rest of the
script: statuses, a missing command, a failed \code{cd}.}

Why are \code{cd} and \code{exit} built in? \code{exit} must end the shell, not a \omterm{def:b2:processes-and-system-calls:tree}{child}. For \code{cd}, the reason is
a theorem about processes.

\begin{definition}[Built-in command]\label{def:b2:processes-and-system-calls:builtin}
A \emph{built-in command}\index{built-in command} is a command that the shell carries out itself, in its own
\omterm{def:b2:processes-and-system-calls:process}{process}, instead of running a program in a \omterm{def:b2:processes-and-system-calls:tree}{child}.
\end{definition}

\begin{proposition}[A child cannot move its parent]\label{prop:b2:processes-and-system-calls:cd-builtin}
A change of working directory made by a \omterm{def:b2:processes-and-system-calls:process}{process} after a \code{fork} has no effect on the working directory of its
\omterm{def:b2:processes-and-system-calls:tree}{parent}. Hence a shell's \code{cd} that changes the shell's own directory must be built in.
\end{proposition}

\begin{proof}
The working directory is part of each \omterm{def:b2:processes-and-system-calls:process}{process}'s state, which \code{fork} copies (\cref{def:b2:processes-and-system-calls:fork}).
\code{chdir} changes the working directory of the \omterm{def:b2:processes-and-system-calls:process}{process} that calls it only; the copy in the \omterm{def:b2:processes-and-system-calls:tree}{parent} is a separate
value, which no call of the \omterm{def:b2:processes-and-system-calls:tree}{child} touches. If \code{cd} were a program, the shell would \omterm{def:b2:processes-and-system-calls:fork}{fork} a \omterm{def:b2:processes-and-system-calls:tree}{child} and the \omterm{def:b2:processes-and-system-calls:tree}{child},
after its \omterm{def:b2:processes-and-system-calls:exec}{exec}, would change its own directory and end, leaving the shell where it was.
\end{proof}

\omcode{code/bachelor-2/13-processes-and-system-calls/listings/cd_child.c}{12}{20}{A \omterm{def:b2:processes-and-system-calls:tree}{child}
moves}

\omoutput{code/bachelor-2/13-processes-and-system-calls/outputs/cd_child.txt}{Its \omterm{def:b2:processes-and-system-calls:tree}{parent} does not.}

The same proposition explains why a script run as a program, \code{./setup.sh}, cannot change the directory of the
shell that started it: the script runs in a \omterm{def:b2:processes-and-system-calls:tree}{child} shell.

\begin{inpractice}[Real shells]
A real shell adds quoting, variables, redirections and pipes (the next chapter's subject), jobs that run in the
background, and dozens of built-ins: \code{export}, \code{source} and \code{umask} must be built in for the reason of
\cref{prop:b2:processes-and-system-calls:cd-builtin}, others such as \code{echo} only for speed. Its heart is still
the loop of \cref{met:b2:processes-and-system-calls:shell-loop}. Processes are also the way servers isolate work:
a web server that \omterm{def:b2:processes-and-system-calls:fork}{forks} a worker per task gets, from the kernel, a worker that cannot corrupt the others' memory.
\end{inpractice}

\section{Exercises}\label{sec:b2:processes-and-system-calls:exercises}

\begin{exercise}[$\star$]\label{exo:b2:processes-and-system-calls:1}
How many lines does this program print? Draw its \omterm{def:b2:processes-and-system-calls:tree}{process tree}.
\omcode{code/bachelor-2/13-processes-and-system-calls/listings/exo_count.c}{5}{11}{Two \omterm{def:b2:processes-and-system-calls:fork}{forks}, one
\code{printf}}
\end{exercise}

\begin{exercise}[$\star$]\label{exo:b2:processes-and-system-calls:2}
This \omterm{def:b2:processes-and-system-calls:tree}{parent} never waits. Explain the output, what remains of each \omterm{def:b2:processes-and-system-calls:tree}{child}, and when it would disappear.
\omcode{code/bachelor-2/13-processes-and-system-calls/listings/exo_nowait.c}{26}{36}{Three children, no
\code{wait}}
\omoutput{code/bachelor-2/13-processes-and-system-calls/outputs/exo_nowait.txt}{Counted from \code{/proc}.}
\end{exercise}

\begin{exercise}[$\star$]\label{exo:b2:processes-and-system-calls:3}
Write \code{runit}, a program that runs the command given on its own command line and exits with that command's
status, reported as a shell would: \code{runit false} exits with 1, \code{runit sh -c 'kill $$'} with 143, and a
missing command with 127.
\end{exercise}

\begin{exercise}[$\star$]\label{exo:b2:processes-and-system-calls:4}
A loop \code{for (int i = 0; i < n; i++) fork();} runs in one \omterm{def:b2:processes-and-system-calls:process}{process}, and nobody waits. How many processes are
created by the loop? How many times is \code{fork} called in all? How many of the final processes have no children?
\end{exercise}

\begin{exercise}[$\star\star$]\label{exo:b2:processes-and-system-calls:5}
In which order are the letters printed? Could another order appear on another run?
\omcode{code/bachelor-2/13-processes-and-system-calls/listings/exo_order.c}{12}{34}{Letters from three
processes}
\end{exercise}

\begin{exercise}[$\star\star$]\label{exo:b2:processes-and-system-calls:6}
This program runs three commands, but a failed \omterm{def:b2:processes-and-system-calls:exec}{exec} is not followed by \code{_exit}. Explain every line of its output.
\omcode{code/bachelor-2/13-processes-and-system-calls/listings/exo_noexit.c}{6}{22}{A missing
\code{_exit}}
\omoutput{code/bachelor-2/13-processes-and-system-calls/outputs/exo_noexit.txt}{One \code{done} too many.}
\end{exercise}

\begin{exercise}[$\star\star$]\label{exo:b2:processes-and-system-calls:7}
Write \code{xargs_lite n command [argument...]}, which reads words from its standard input and runs the command with
its arguments followed by at most $n$ of the words at a time, until the words run out, and exits with the largest
status it met. \code{printf 'a b c d e' | xargs_lite 2 echo X} prints \code{X a b}, \code{X c d}, \code{X e}.
\end{exercise}

\begin{exercise}[$\star\star$]\label{exo:b2:processes-and-system-calls:8}
Prove \cref{prop:b2:processes-and-system-calls:fork-count} by induction on $n$. How many processes print in
\code{fork_tree.c} with $n$ \omterm{def:b2:processes-and-system-calls:fork}{forks}, and how many of them have depth $d$?
\end{exercise}

\begin{exercise}[$\star\star$]\label{exo:b2:processes-and-system-calls:9}
Write a built-in \code{cd} that also understands \code{cd -}, ``back to the previous directory'', and prints the
directory it returns to. Why must the previous directory be stored in the shell, and when must it be updated?
\end{exercise}

\begin{exercise}[$\star\star\star$]\label{exo:b2:processes-and-system-calls:10}
Add \code{&&} and \code{||} to \code{minish}: in \code{a && b}, \code{b} runs only if \code{a} succeeded; in
\code{a || b}, only if it failed; the operators chain from left to right with equal precedence, as in the usual
shells. Write \code{run_list(argv, run)}, which takes the words of a line and a function running one command, and
returns the status of the last command run. What does \code{false && x || true} run?
\end{exercise}

\begin{exercise}[$\star\star\star$]\label{exo:b2:processes-and-system-calls:11}
\code{lines.c} prints $n$ numbered lines of 10 bytes, by \code{printf}, with \code{printf}'s buffer or without it.
Count its \code{write} \omterm{def:b2:processes-and-system-calls:syscall}{system calls} for $n = 1000$ in both modes, explain the two counts, and predict them for
$n = 10\,000$. What would a terminal change?
\omoutput{code/bachelor-2/13-processes-and-system-calls/outputs/count_writes.txt}{Each \code{write} is one line of
\code{strace}'s output.}
\end{exercise}

\begin{exercise}[$\star\star\star$]\label{exo:b2:processes-and-system-calls:12}
Write \code{ptree}, a program that starts two children, the first with two children of its own, gives them names, and
prints its own \omterm{def:b2:processes-and-system-calls:tree}{process tree} from \code{/proc}, indented, as \code{pstree} would, before stopping them all. On Linux,
\code{/proc/<pid>/task/<pid>/children} lists the PIDs of a \omterm{def:b2:processes-and-system-calls:process}{process}'s children and \code{/proc/<pid>/comm} holds its
name. Make the output the same on every run.
\end{exercise}

\section{Problem: Building the tiny shell}\label{sec:b2:processes-and-system-calls:problem}

\begin{problem}[{Weekend problem --- Building the tiny shell}]\label{pb:b2:processes-and-system-calls:1}
The problem rebuilds \code{minish} step by step, then counts the \omterm{def:b2:processes-and-system-calls:syscall}{system calls} it makes to run one command.

\medskip\noindent\textbf{Part I --- Parsing.}\par\nopagebreak
\begin{enumerate}
\item What does \code{split} produce for a line made of two spaces, \code{echo}, three spaces, \code{two}, a tab,
\code{word} and a newline? Where do the strings of the vector live?
\item Why must the vector end with a null pointer?
\item What happens to the line \code{echo $? $?} after \code{false}? To \code{echo x$?}?
\item \code{strtok} keeps its position in a hidden static variable. Why is that harmless in \code{minish} and not in
a recursive function?
\item A line of 2000 characters is typed. What does \code{fgets} with a buffer of 1024 do, and what does the shell
then run?
\item Give a command line that \code{minish} cuts differently from a real shell.
\end{enumerate}

\medskip\noindent\textbf{Part II --- Fork, exec, wait.}\par\nopagebreak
\begin{enumerate}[resume, beginpenalty=10000]
\item In \code{run}, which lines run in the \omterm{def:b2:processes-and-system-calls:tree}{parent}, which in the \omterm{def:b2:processes-and-system-calls:tree}{child}, and which in both?
\item Why does the shell call \code{fflush} before forking? Describe the output of \code{minish} on a script
without it.
\item What does the \omterm{def:b2:processes-and-system-calls:tree}{child} do if \code{execvp} succeeds? What does it do if it fails, and why \code{_exit} rather than
\code{return}?
\item Where does \code{execvp} look for \code{echo}? What changes if the line is \code{/usr/bin/echo hi}?
\item Why does the shell wait before printing the next prompt? What would a script's output look like if it did not?
\item How many processes exist, at most, while \code{minish} runs the script of the hook?
\end{enumerate}

\medskip\noindent\textbf{Part III --- Statuses and built-ins.}\par\nopagebreak
\begin{enumerate}[resume, beginpenalty=10000]
\item Explain the four statuses printed by \code{echo $?} in the hook's script.
\item What status does \code{run} return for a command killed by signal 9? Why add 128?
\item What does \code{minish} return when its input ends? When it reads \code{exit}? \code{exit 300}?
\item Why is \code{exit} built in? Use \cref{prop:b2:processes-and-system-calls:cd-builtin} to say why \code{cd} is.
\item \code{pwd} is not built into \code{minish}, yet it prints the directory \code{cd} chose. Why?
\item Name one command that could be built in or not, with the same result, and one that could not.
\end{enumerate}

\medskip\noindent\textbf{Part IV --- The system calls of one command.}\par\nopagebreak
\omoutput{code/bachelor-2/13-processes-and-system-calls/outputs/budget.txt}{The calls of \code{minish} running
\code{echo hi}, counted by \code{strace -f -c}, with the calls that failed; then two
baselines.}\label{out:b2:processes-and-system-calls:budget}
\begin{enumerate}[resume, beginpenalty=10000]
\item How many \omterm{def:b2:processes-and-system-calls:syscall}{system calls} does running \code{echo hi} through \code{minish} add to the shell's own 41? How many of
them does \code{echo} make by itself?
\item Of the 6 calls left, the trace shows that 5 are made by the shell and 1 by the \omterm{def:b2:processes-and-system-calls:tree}{child} before its \omterm{def:b2:processes-and-system-calls:exec}{exec}. Find
the shell's 5 from the code of \code{run} and of the loop. Under which names do \code{fork} and \code{waitpid}
appear in the table?
\item Why are there 2 \code{execve} calls? Why would there be more with a longer \code{PATH}?
\item \code{echo} alone makes 115 calls to print 3 bytes. Most of them are \code{openat}, \code{mmap},
\code{newfstatat}, \code{close}. What is it doing before it can print?
\item Using \cref{out:b2:processes-and-system-calls:cost}, estimate the time of these 162 calls, if each cost what
\code{getppid} costs. Is that the time of the command?
\item State the system-call budget of one external command.
\end{enumerate}
\end{problem}
